\appendices

\section{Correctness of the baseline}
\label{app:baseline}
The proofs in these appendices fix a strict total order $\prec$ on the items, in which
every class lists its members and joins a member only with the members after it. The
baseline orders items by identifier. The optimized variant orders them by increasing
support and breaks ties by identifier (Section~\ref{sec:optimized}). No proof uses
anything about $\prec$ beyond it being a strict total order, so every result holds for
both variants. We write $x \preceq y$ when $x \prec y$ or $x = y$. For a nonempty itemset
$X$, $\ell(X)$ is its last item under $\prec$, and every item comes after
$\ell(\emptyset)$.

A class with prefix $P$ is a list of members $(b, T_b)$, where $b$ is an item and $T_b$ a
tid-list. The search \texttt{bottom\_up!} visits the members of a class in order. For each
member $a$ it records $P \cup \{a\}$ with support $|T_a|$, intersects $T_a$ with $T_b$ for
every member $b$ after $a$, keeps each result of size at least $\sigma_{\min}$ as a
member of the child class with prefix $P \cup \{a\}$, and searches that child class before
it moves on. The root class has prefix $\emptyset$ and holds every frequent item with its
tidset.

\begin{lemma}[Invariants]
\label{lem:invariants}
Every class with prefix $P$ that the search receives satisfies: (i) each tid-list is
increasing and has no duplicates; (ii) the members are in increasing order under $\prec$,
and every member item comes after $\ell(P)$; (iii) every member satisfies
$|T_b| \ge \sigma_{\min}$; (iv) every member satisfies $T_b = t(P \cup \{b\})$.
\end{lemma}
\begin{IEEEproof}
By induction on the depth of the class. The root class satisfies (i) because
\texttt{to\_class} appends tids in the order it reads the transactions, (ii) because the
root is sorted and $P = \emptyset$, (iii) because it keeps only frequent items, and (iv)
by the definition of $t(\{b\})$. Let a class with prefix $P$ satisfy (i) to (iv), and let
its member $a$ produce the child class with prefix $P' = P \cup \{a\}$. The merge of two
increasing lists without duplicates is increasing without duplicates, which gives (i).
The child keeps only members $b$ after $a$, in their original order, and $\ell(P') = a$,
which gives (ii). The size check before a member is kept gives (iii). For (iv),
identity~\eqref{eq:intersect} gives
$T_a \cap T_b = t(P \cup \{a\}) \cap t(P \cup \{b\}) = t(P' \cup \{b\})$.
\end{IEEEproof}

\begin{lemma}[Exact support]
\label{lem:exact}
Every recorded itemset is frequent and is recorded with its support.
\end{lemma}
\begin{IEEEproof}
The search records $P \cup \{a\}$ with $|T_a|$. By
Lemma~\ref{lem:invariants}(iv) and~\eqref{eq:support}, $|T_a| = |t(P \cup \{a\})| =
\sigma(P \cup \{a\})$, and by (iii) this is at least $\sigma_{\min}$.
\end{IEEEproof}

\begin{lemma}[No duplicates]
\label{lem:unique}
No itemset is recorded twice.
\end{lemma}
\begin{IEEEproof}
First, each prefix $P$ gives at most one class. The root is built once. A class with a
nonempty prefix $P$ is built only while the class with prefix $P \setminus \{\ell(P)\}$
visits its member $\ell(P)$. By induction on $|P|$ that parent class is built at most
once, and it visits each member once, because its members are distinct by
Lemma~\ref{lem:invariants}(ii). Second, when the class with prefix $P$ records
$X = P \cup \{a\}$, Lemma~\ref{lem:invariants}(ii) gives $a = \ell(X)$ and therefore
$P = X \setminus \{\ell(X)\}$. Both the class and the member that record $X$ are thus
determined by $X$, and each records once.
\end{IEEEproof}

\begin{lemma}[Anti-monotone support]
\label{lem:antimonotone}
If $X \subseteq Y$, then $t(Y) \subseteq t(X)$ and $\sigma(Y) \le \sigma(X)$. Every subset
of a frequent itemset is therefore frequent.
\end{lemma}
\begin{IEEEproof}
A transaction $T_i$ with $Y \subseteq T_i$ also satisfies $X \subseteq T_i$, so by
\eqref{eq:support} every tid in $t(Y)$ is in $t(X)$. Taking sizes gives the inequality.
\end{IEEEproof}

\begin{lemma}[Completeness]
\label{lem:complete}
Every frequent itemset is recorded.
\end{lemma}
\begin{IEEEproof}
Let $X = \{x_1 \prec \dots \prec x_k\}$ be frequent and let $P_j = \{x_1, \dots, x_j\}$. We
show by induction on $j = 1, \dots, k$ that the search builds the class with prefix
$P_{j-1}$, and that this class contains every item $y$ after $\ell(P_{j-1})$ for which
$P_{j-1} \cup \{y\}$ is frequent. For $j = 1$, the prefix is empty and the root holds
every frequent item. Assume the claim for $j < k$. The itemset $P_j \subseteq X$ is
frequent by Lemma~\ref{lem:antimonotone}, and $x_j$ comes after $\ell(P_{j-1}) =
x_{j-1}$, so $x_j$ is a member of the class with prefix $P_{j-1}$. Visiting $x_j$ builds
the class with prefix $P_j$. Now take $y$ after $x_j$ with $P_j \cup \{y\}$ frequent.
Lemma~\ref{lem:antimonotone} makes $P_{j-1} \cup \{y\}$ frequent, so $y$ is a member of
the parent class after $x_j$. The search intersects their tid-lists, the result is
$t(P_j \cup \{y\})$ by Lemma~\ref{lem:invariants}(iv), its size is at least
$\sigma_{\min}$, and the child keeps $y$. This proves the claim for $j + 1$. For
$j = k$, the class with prefix $P_{k-1}$ contains $x_k$, and visiting $x_k$ records $X$.
\end{IEEEproof}

\begin{theorem}
\label{thm:baseline}
The baseline records exactly the frequent itemsets, each once, with its support.
\end{theorem}
\begin{IEEEproof}
Lemmas~\ref{lem:exact}, \ref{lem:unique}, and~\ref{lem:complete}.
\end{IEEEproof}

\section{Correctness of perfect-extension pruning}
\label{app:perfect}
With perfect-extension pruning, the search no longer records one itemset per member.
Each search node stands for a family of itemsets. A node is a triple $(C, E, M)$ of
pairwise disjoint item sets: the core $C$, the items joined so far; the perfect set $E$;
and the class $M$, the members still to join, in order under $\prec$, each with the
tidset $t(C \cup \{b\})$. The root is $(\emptyset, \emptyset, F)$, where $F$ is the set of
frequent items. For a member $a \in M$, the node $(C, E, M)$ has a child $(C', E', M')$
with
\begin{align}
  C' &= C \cup \{a\}, \nonumber\\
  E' &= E \cup \{\, b \in M : a \prec b,\ \sigma(C' \cup \{b\}) = \sigma(C') \,\},
    \label{eq:perfect-set}\\
  M' &= \{\, b \in M : a \prec b,\ \sigma_{\min} \le \sigma(C' \cup \{b\}) < \sigma(C') \,\}.
    \nonumber
\end{align}
Every node except the root emits the itemset $C \cup S$ with support $\sigma(C)$ for
every $S \subseteq E$, which gives $2^{|E|}$ itemsets. The search computes each
$\sigma(C' \cup \{b\})$ by one intersection, as in Appendix~\ref{app:baseline}, so
Lemma~\ref{lem:invariants}(iv) carries over and every tidset that a member carries is
exact. The three sets stay disjoint: $a \in M$ is not in $E$, the new items of $E'$ and
all of $M'$ come after $a = \ell(C')$ and so lie outside $C'$, and $M'$ excludes the
items that enter $E'$. Every core item is also a former member of $M$, so
$\ell(C) \prec \ell(C')$ and the core items arrive in increasing order along a path.

\begin{lemma}[Exact support]
\label{lem:perfect-exact}
For every node $(C, E, M)$ other than the root and every $S \subseteq E$,
$t(C \cup S) = t(C)$ and $\sigma(C) \ge \sigma_{\min}$. Every emitted itemset is
therefore frequent and is emitted with its support.
\end{lemma}
\begin{IEEEproof}
By induction on depth. The identity $t(C \cup S) = t(C)$ also holds at the root,
where $E = \emptyset$. Let it hold at $(C, E, M)$, and let $(C', E', M')$ be the child of member
$a$. Split $S \subseteq E'$ into $S_0 = S \cap E$ and $S_1 = S \setminus E$. Then
\begin{equation*}
  t(C' \cup S_0) = t(C \cup S_0) \cap t(\{a\}) = t(C) \cap t(\{a\}) = t(C').
\end{equation*}
Each $b \in S_1$ satisfies $\sigma(C' \cup \{b\}) = \sigma(C')$ by~\eqref{eq:perfect-set},
and $t(C' \cup \{b\}) \subseteq t(C')$ by Lemma~\ref{lem:antimonotone}. A subset with the
same size is the whole set, so $t(C' \cup \{b\}) = t(C')$ and $t(C') \subseteq t(\{b\})$.
Intersecting $t(C' \cup S_0) = t(C')$ with $t(\{b\})$ for every $b \in S_1$ therefore
leaves $t(C')$, which proves $t(C' \cup S) = t(C')$. Finally $a \in M$, and every member
$b$ of a class satisfies $\sigma(C \cup \{b\}) \ge \sigma_{\min}$: the root holds frequent
items, and $M'$ keeps only such members. Hence $\sigma(C') \ge \sigma_{\min}$.
\end{IEEEproof}

\begin{lemma}[Class membership]
\label{lem:perfect-members}
For every node $(C, E, M)$,
$M = \{\, b : \ell(C) \prec b,\ b \notin E,\ \sigma(C \cup \{b\}) \ge \sigma_{\min} \,\}$.
\end{lemma}
\begin{IEEEproof}
At the root both sides are $F$. Let the claim hold at $(C, E, M)$ and let
$(C', E', M')$ be the child of $a$, so $\ell(C') = a$. If $b \in M'$, then $a \prec b$,
$C' \cup \{b\}$ is frequent, $b \notin E$ because $b \in M$, and $b$ is not a new item of
$E'$ because its support is below $\sigma(C')$, so $b$ is in the right-hand side.
Conversely, take $b$ with $a \prec b$, $b \notin E'$, and $C' \cup \{b\}$ frequent. Then
$C \cup \{b\}$ is frequent by Lemma~\ref{lem:antimonotone}, $\ell(C) \prec a \prec b$,
and $b \notin E \subseteq E'$, so $b \in M$ by the induction hypothesis. Since $b$ did
not enter $E'$, $\sigma(C' \cup \{b\}) < \sigma(C')$, and $b \in M'$.
\end{IEEEproof}

For a frequent itemset $Z \neq \emptyset$, the \emph{walk} of $Z$ starts at the root and,
at each node $(C, E, M)$ with $C \subseteq Z$, looks at the remainder
$R = Z \setminus (C \cup E)$. It stops if $R = \emptyset$, and otherwise moves to the
child of the first item of $R$ under $\prec$. The walk is deterministic: $Z$ alone decides
each step.

\begin{lemma}[Coverage]
\label{lem:perfect-cover}
The walk of every nonempty frequent itemset $Z$ ends at a node other than the root, and
that node emits $Z$.
\end{lemma}
\begin{IEEEproof}
We show that the walk can always take its step and keeps two properties: $C \subseteq Z$,
and every item of $R$ comes after $\ell(C)$. Both hold at the root. At a node where
$R \neq \emptyset$, let $a$ be the first item of $R$. It comes after $\ell(C)$, it is not
in $E$, and $C \cup \{a\} \subseteq Z$ is frequent by Lemma~\ref{lem:antimonotone}, so
$a \in M$ by Lemma~\ref{lem:perfect-members} and the child of $a$ exists. At the child,
$C' = C \cup \{a\} \subseteq Z$, and its remainder is a subset of $R \setminus \{a\}$,
whose items all come after $a = \ell(C')$. The remainder loses at least $a$ at every
step, so the walk stops. It does not stop at the root, where $R = Z \neq \emptyset$. At
the final node, $R = \emptyset$ gives $Z = C \cup (Z \cap E)$ with $Z \cap E \subseteq E$,
so the node emits $Z$.
\end{IEEEproof}

\begin{lemma}[Uniqueness]
\label{lem:perfect-unique}
No itemset is emitted twice.
\end{lemma}
\begin{IEEEproof}
Within one node, two different subsets $S \neq S'$ of $E$ give $C \cup S \neq C \cup S'$,
because $E$ is disjoint from $C$. It remains to show that only one node emits a given $Z$.
Let a node $N_k$ emit $Z = C_k \cup S$ with $S \subseteq E_k$, and let
$N_0, N_1, \dots, N_k$ be the path from the root, where $N_{j+1}$ is the child of member
$a_{j+1}$ of $N_j = (C_j, E_j, M_j)$. We show that this path is the walk of $Z$. Since
each node has a single path from the root and the walk is deterministic, $N_k$ is then
the only node that emits $Z$.

Fix $j < k$ and let $R_j = Z \setminus (C_j \cup E_j)$. The item $a_{j+1}$ is in $Z$
because $a_{j+1} \in C_k$, it is not in $C_j$ because it comes after $\ell(C_j)$, and it
is not in $E_j$ because $a_{j+1} \in M_j$. So $a_{j+1} \in R_j$, and the walk does not
stop at $N_j$. Suppose some $r \in R_j$ comes before $a_{j+1}$. If $r \in C_k$, then
$r = a_i$ for some $i > j$, but core items arrive in increasing order, so
$a_{j+1} \preceq a_i = r$, a contradiction. Otherwise $r \in S \subseteq E_k$ and
$r \notin E_j$, so $r$ entered the perfect set at some step $i > j$. By
\eqref{eq:perfect-set}, the items that enter at step $i$ come after $a_i$, and
$a_{j+1} \preceq a_i$, again a contradiction. Hence $a_{j+1}$ is the first item of
$R_j$, which is the step the walk takes. At $N_k$ the remainder is
$Z \setminus (C_k \cup E_k) = \emptyset$, so the walk stops there.
\end{IEEEproof}

\begin{theorem}
\label{thm:perfect}
The search with perfect-extension pruning emits exactly the frequent itemsets, each once,
with its support.
\end{theorem}
\begin{IEEEproof}
Lemmas~\ref{lem:perfect-exact}, \ref{lem:perfect-cover}, and~\ref{lem:perfect-unique}.
\end{IEEEproof}

\section{Exactness of the skipped and stopped intersections}
\label{app:skips}
The optimized variant changes how a tidset is stored and skips or stops some
intersections. None of these changes alters a node.

\emph{Representations.} A bitset sets bit $i$ exactly when tid $i$ is in the tidset and
leaves the unused high bits of its last word at zero, so the word-by-word AND of two
bitsets sets exactly the bits of the common tids, and its population count is the size of
the intersection. Converting a bitset to a tid list lists exactly its set bits, so a
member keeps the same tidset in a new form.

\begin{lemma}
\label{lem:skip}
Skipping the intersection of $a$ with a later member $b$ whose result is known to satisfy
$\sigma(C' \cup \{b\}) < \sigma_{\min}$ leaves the child $(C', E', M')$ unchanged.
\end{lemma}
\begin{IEEEproof}
By~\eqref{eq:perfect-set}, $b$ enters $E'$ only if $\sigma(C' \cup \{b\}) = \sigma(C')$,
which is at least $\sigma_{\min}$ by Lemma~\ref{lem:perfect-exact}, and it enters $M'$
only if $\sigma(C' \cup \{b\}) \ge \sigma_{\min}$. An infrequent result therefore places
$b$ in neither set, whether the intersection runs or not.
\end{IEEEproof}

Each of the three rules below skips or stops an intersection only when an upper bound on
its support is below $\sigma_{\min}$, so Lemma~\ref{lem:skip} applies to all of them.

\emph{Pair counts at the root.} At the root $C = \emptyset$ and $C' = \{a\}$, and the
count of the pair $(a, b)$ is the number of transactions that contain both items, which
is $\sigma(\{a, b\})$ by~\eqref{eq:support}. The count is exact, so a pair below
$\sigma_{\min}$ is infrequent.

\emph{Early stop in a tid-list merge.} Suppose a merge has found $f$ common tids and has
$r_a$ and $r_b$ tids left in the two lists. Every further common tid is among the
remaining tids of both lists, so the final size is at most $f + \min(r_a, r_b)$. The
merge stops only when this bound is below $\sigma_{\min}$. It never stops on reaching
$\sigma_{\min}$, so a frequent result carries its whole tidset and its exact support,
which both the perfect-extension test and the child class need.

\emph{Block bounds.} Split the tids into disjoint blocks. Each common tid of $X$ and $Y$
lies in exactly one block $j$, and the number of common tids in block $j$ is at most
$\min(c_{X,j}, c_{Y,j})$, the smaller of the two members' counts there. Summing over the
blocks gives~\eqref{eq:bound}, so a sum below $\sigma_{\min}$ proves the intersection
infrequent.

\begin{theorem}
The optimized variant emits exactly the frequent itemsets, each once, with its support.
\end{theorem}
\begin{IEEEproof}
By Lemma~\ref{lem:skip} and the three bounds above, the optimized search builds the same
nodes as the search of Appendix~\ref{app:perfect} under its item order. That search is
correct by Theorem~\ref{thm:perfect}.
\end{IEEEproof}
